\section{Individual Contribution}\label{sec:individual-contribution}

This dissertation is an individual thesis on Component 01 of the NeuroShell
four-component research chain. The candidate's individual contribution is the
design, implementation and evaluation of the Intent Recognition Engine, and
comprises: (1) the architecture of the guarded NLU gateway, in which a small
offline language model acts as a proposition generator and deterministic stages
validate structure, target syntax and authorisation before an intent is
emitted; (2) the security countermeasures layered around the core---including
the authorisation checks, the adversarial-utterance detection, and the
self-consistency and hallucination-validation gate---that enforce the safety
invariant that every accepted intent is schema-valid, authorised and in scope;
(3) the authoring of the annotated evaluation data set covering well-formed and
deliberately flawed prompt categories; (4) the evaluation harness and the
per-layer ablation protocol used to quantify the pipeline's value and each
layer's contribution; and (5) the empirical comparison of two offline
language-model backends under the identical guardrail stack.

The other components of the NeuroShell chain---the dynamic planner, the
adaptive executor, and the vulnerability-analysis engine---provide the context
in which the IRE operates and are described only to the extent needed to define
the IRE's interface contract. The candidate takes responsibility for the
correctness of the IRE implementation and for every result reported in this
dissertation.